\paragraph{Which part of a history depends on the observer?}
Fix the matrix rank, a material clock sector and a network of physical contacts. At contact \(e\), let \(L_{e,z}=\mathbb W_e K_{e,z}\) for receipt \(z\); use \(\mathbb W_e=\id\) if no separate routing occurs. The clock orders the contacts. The factors may depend on receipts already available at that port. If \(R_O({\bf z})\) is the part of the record that reaches observer \(O\), its operation on the common preparation \(\rho\) is
\begin{equation}
 \mathcal I^O_r(\rho)=\sum_{R_O({\bf z})=r} L_{\bf z}\rho L_{\bf z}^{\dagger},
 \qquad L_{\bf z}=L_{e_s,z_s}\cdots L_{e_1,z_1}.
 \label{eq:matrixhistory}
\end{equation}
The trace gives the Born weight and normalization gives the conditional state. On this fixed network every fine history occurs once in the sum over \(r\). All observers therefore have the same total channel when the outcomes are discarded. This remains true with adaptive contacts if the operations are complete after each reachable sequence of receipts. The composition of instruments is standard quantum-network theory \cite{networks}.

CAT makes a further physical prescription. Each observer uses the material port reached by that observer's signals, the charge recorded there at the deadline, and the output modes allowed by that charge. A recorded charge controls a subsequent operation even if no one later reads the register. Section~\ref{sec:joint} follows one qubit from a common preparation. On a joint receipt of weight \(0.196608\), A's port admits an inward copy of the \emph{code space}; B's port has zero inward dimension and admits a surface carrier. Each map acts on the one unknown input in its own description. In a common causal region a delivered copy of the same classical receipt must give the same probabilities. The supplement specifies the maps, the accessible observables and the outcome-wise clock-order condition.

\paragraph{Matrix size and the scope of the proposal.}
The ten labels \(a\) in Eq.~\eqref{eq:ikkt} count matrix components; \(N\) is the size of each matrix. CAT fixes neither a highest \(N\) nor a universal number of gravitational contacts. For any admissible preparation, Eq.~\eqref{eq:matrixhistory} uses the specified scattering blocks with their full coupling dependence. This is the sense in which the representation is nonperturbative. A longer network allows further receipts, and a continuous outcome requires a normalized instrument on measurable sets. These possibilities are distinct from taking \(N\) to infinity.

The original BFSS and IKKT proposals both give a physical role to large \(N\) \cite{bfss,ikkt}. CAT adds irreversible contacts and observer-specific output spaces, so the ordinary matrix limit by itself would leave a further question: which receipts and observables at two ranks represent the same physical preparation? A completion would have to identify those observables and obtain consistent probabilities, whether by a controlled limit or by a physical rule for weighting sectors. Section~\ref{sec:geometry} removes prescribed coordinate time on a chosen material-clock sector. Neither that construction nor the finite examples supply the general measure.

